\documentclass{article}
\usepackage[utf8]{inputenc}
\usepackage[T1]{fontenc}
\usepackage{geometry}
\geometry{
    a4paper,
    left=2.5cm,
    right=2.5cm,
    top=2.5cm,
    bottom=2.5cm
}
\usepackage{graphicx}
\usepackage{float}
\usepackage{hyperref}
\usepackage[english]{babel}
\usepackage{amsmath}
\usepackage{amssymb}
\usepackage{array}
\usepackage{booktabs}
\usepackage{caption}
\usepackage{subcaption}
\usepackage{enumitem}
\usepackage{fancyhdr}
\usepackage{titlesec}
\usepackage{url}
\usepackage{textcomp}
\usepackage{xcolor}
\usepackage{listings}
\usepackage{multirow}
\usepackage{longtable}
\usepackage{tabularx}
\usepackage{times}

% Set up page style
\pagestyle{fancy}
\fancyhf{}
\rfoot{\thepage}
\renewcommand{\headrulewidth}{0pt}

% Title formatting
\titleformat{\chapter}[display]{\normalfont\bfseries}{\chaptertitle\thechapter}{0pt}{\Huge}
\titlespacing*{\chapter}{0pt}{50pt}{40pt}

% Document info
\title{Template Document}
\author{Academic Support Agent}
\date{\today}

\begin{document}
\maketitle

\begin{frame}
  \titlepage
\end{frame}

\section{Secci n}
\subsection{Subsecci n}
\begin{frame}
  \frametitle{T tulo diapositiva}
  Esperanza $\esperanza{X}$. probabilidad $\proba{X}$ \begin{theorem} teorema
      
  \end{theorem}
  
  \begin{proof} Demostraci n \end{proof}

De acuerdo a \cite{dirac}, o en estilo natbib \citap{dirac}.
\end{frame}

\section{Gr ficas}
\subsection{tikz}

\begin{frame}[fragile]{Como usar tikz}

\begin{tikzcd}
A \dar{g} \rar{f} & B \dar{h} \\
C \rar{i} & D
\end{tikzcd}
\\
Para que tikz funcione se deben definir los frames as : begin{frame}[fragile]{Como usar tikz}.

\end{frame}

\begin{frame}[fragile]{Otro ejemplo}
    \begin{center}
    

\begin{tikzcd}
      & Y \arrow[d, "f" description, dashed] \arrow[rd, "f_2"] \arrow[ld, "f_1"'] &       \\
X_{1} & X_{1} \times X_{2} \arrow[r, "\pi_2"'] \arrow[l, "\pi_2"]                 & X_{2}
\end{tikzcd}

\end{center}
\end{frame}

\section{Bibliograf a}
\begin{frame}{Bibliograf a}
\nocite{*}
\printbibliography
\end{frame}
\end{document}